% SPDX-License-Identifier: LGPL-3.0-only
\chapter{Numerical Integration}
\label{ch:integration}

\section{The Midpoint Rule}
For a continuous function \(f\), approximate its integral by sampling
the midpoint of each subinterval~\cite{calculus}.
Appendix~\ref{app:notation} collects the notation.

\begin{definition}[Composite midpoint rule]
For \(n\geq 1\) and \(h=(b-a)/n\), define
\begin{equation}
  M_n(f) = h\sum_{k=0}^{n-1}
    f\left(a+\left(k+\frac{1}{2}\right)h\right).
  \label{eq:midpoint}
\end{equation}
\end{definition}

\begin{theorem}[Exactness for affine functions]
\label{thm:affine}
If \(f(x)=cx+d\), then \(M_n(f)=\int_a^b f(x)\,\mathrm{d}x\).
\end{theorem}
\begin{proof}
On each subinterval, the average of an affine function equals its
value at the midpoint. Multiplying by \(h\) gives the exact integral
on that subinterval. Summing proves the result.
\end{proof}

\begin{example}
For \(f(x)=x\) on \([0,1]\), Equation~\eqref{eq:midpoint} gives
\(M_n(f)=1/2\) for every positive integer \(n\).
\end{example}

\begin{exercise}
Apply the midpoint rule to \(f(x)=x^2\) on \([0,1]\) with \(n=2\).
Compare the result with the exact integral \(1/3\).
\end{exercise}
